\section{Relate counts, symbol positions and arriving code bits}
The fixed alphabet gives three different numbers for each populated length. The count
states how many symbols have that length. The exclusive offset states where those
symbols begin in the ordered array. The first code states their smallest code value.
Do not use one of these numbers as a substitute for another.

\Needspace{9\baselineskip}
\begin{longtable}{p{24mm}p{30mm}p{36mm}p{47mm}}
  \toprule Code length & Symbol count & Ordered-array start & First canonical code \\\midrule
  7                    & 24           & 0                   & 0                    \\
  8                    & 152          & 24                  & 48                   \\
  9                    & 112          & 176                 & 400                  \\\bottomrule
\end{longtable}
After failing to match at length seven, the decoder advances first code to
$(0+24)\times2=48$ and ordered-array start to 24. Arrivals for literal zero complete
code 48 at length eight. Difference zero selects array index $24+0$, whose stored
symbol is zero. This trace explains why neither code 48 nor array index 24 should be
written as the output byte.

Scratch placement offsets are used while constructing a table and then reused.
The decoder reconstructs the current first code and array start from the immutable
counts of that completed table. It does not depend on the scratch cursors retaining
their original prefix values after sorting the symbols.
